\documentclass[lettersize,journal]{IEEEtran}
\usepackage{iftex}
\ifPDFTeX\else
  \PackageError{bare_jrnl_new_sample4}{Use pdfLaTeX, not XeLaTeX or LuaLaTeX}{Select the LaTeX Workshop recipe: pdflatex-shellescape -> bibtex -> pdflatex x2.}
\fi
\usepackage{amsmath,amsfonts}
\usepackage{algorithmic}
\usepackage{algorithm}
\usepackage{array}
\usepackage[caption=false,font=normalsize,labelfont=sf,textfont=sf]{subfig}
\usepackage{textcomp}
\usepackage{stfloats}
\usepackage{url}
\usepackage{verbatim}
\usepackage{graphicx}
\usepackage{svg}
\usepackage{cite}
\DeclareUnicodeCharacter{03BB}{$\lambda$}
\DeclareUnicodeCharacter{00D7}{$\times$}
\def\BibTeX{{\rm B\kern-.05em{\sc i\kern-.025em b}\kern-.08em
    T\kern-.1667em\lower.7ex\hbox{E}\kern-.125emX}}
\usepackage{balance}
\hyphenation{op-tical net-works semi-conduc-tor IEEE-Xplore}
% updated with editorial comments 8/9/2021

\begin{document}

\title{Slim-Mamba: An Efficient and Reconfigurable Mamba-Based Accelerator for Massive MIMO Indoor Localization}

\author{IEEE Publication Technology,~\IEEEmembership{Staff,~IEEE,}
        % <-this % stops a space
\thanks{This paper was produced by the IEEE Publication Technology Group. They are in Piscataway, NJ.}% <-this % stops a space
\thanks{Manuscript received April 19, 2021; revised August 16, 2021.}}

% The paper headers
\markboth{Journal of \LaTeX\ Class Files,~Vol.~14, No.~8, August~2021}%
{Shell \MakeLowercase{\textit{et al.}}: A Sample Article Using IEEEtran.cls for IEEE Journals}

% \IEEEpubid{0000--0000/00\$00.00~\copyright~2021 IEEE}
% Remember, if you use this you must call \IEEEpubidadjcol in the second
% column for its text to clear the IEEEpubid mark.

\maketitle

\begin{abstract}
Massive MIMO radio measurements provide rich spatial and delay-domain information
for indoor localization, but their high-dimensional channel features make
real-time neural inference challenging on resource-constrained hardware.
This paper presents Slim-Mamba, an efficient and reconfigurable Mamba-based
accelerator for Massive MIMO indoor localization. At the algorithm level,
Slim-Mamba simplifies the original Mamba-style selective state-space block into
a lightweight gated recurrent update, preserving temporal state modelling and
output gating while reducing the complexity of the selective-scan computation.
The deployed model contains 816k trainable parameters and reduces the model size
to 3.40\,MB while maintaining competitive localization accuracy. At the
architecture level, the accelerator supports reconfigurable inter-block execution,
MAC-fabric scheduling, and quantization. The Slim-Mamba
blocks can operate with either streaming data transfer or preload execution,
while a shared $4\times4\times4$ MAC fabric is reused across multiple projection operators
through configurable scheduling. In addition, a unified requantization engine
enables several modes for scaling, rounding, and saturation. A fine-grained pipeline further overlaps
projection, nonlinear activation, state update, and output gating, reducing
latency by 23.8\%. Combined with the reconfigurable architecture, the design
achieves both efficient resource utilization and lower inference delay.
Implemented on a Zynq MPSoC FPGA, the accelerator achieves a mean localization
error of 0.14103\,m, close to the float result of 0.13500\,m, and reaches
471\,frames/s at 100\,MHz. The design uses 22.21\% LUTs, 14.11\% FFs, 30.59\% BRAMs,
and 19.68\% DSPs, while the programmable logic (PL) design consumes only 0.94\,W,
demonstrating efficient real-time deployment of Mamba-style localization models
on FPGA.
\end{abstract}

\begin{IEEEkeywords}
Article submission, IEEE, IEEEtran, journal, \LaTeX, paper, template, typesetting.
\end{IEEEkeywords}

\section{Introduction}
\IEEEPARstart{I}{ndoor} localization is becoming an important function in future wireless systems,
where communication, sensing, and positioning are expected to be more tightly
integrated. In the context of 6G, localization has been identified as both a key
service and an enabling technology for location-aware applications such as robotic
navigation, emergency healthcare, smart environments, and intelligent
transportation~\cite{Trevlakis2023}. However, accurate indoor positioning
remains challenging because indoor radio propagation is affected by multipath
effects, noise, and the surrounding radio environment
~\cite{Zafari2019Indoor}.

Massive multiple-input multiple-output (MIMO) systems provide a promising
measurement platform for radio-based indoor localization. The LuViRA dataset used in this paper 
provides synchronized vision, radio, audio, and motion-capture measurements for indoor 
localization research~\cite{LuViRA2024,LuViRAValidation2024}. In its radio modality, the
channel response is measured using the LuMaMi Massive MIMO testbed
~\cite{LuMaMi2017}, enabling the collection of rich Massive MIMO channel information that can be 
used for learning-based localization, making LuViRA a suitable dataset for evaluating
radio-based localization models under practical indoor propagation conditions. 
Previous FCNN-based Massive MIMO localization studies used spatial covariance
matrices and truncated CIRs as fingerprints, and reached centimetre-level
positioning accuracy~\cite{TianICC2023}. However, the resulting FCNN models
require network sizes ranging from hundreds of megabytes to more than 1~GB. The extension 
in~\cite{Tian2024} further confirms that increasing network size is a central challenge when
ML-based localization is applied to Massive MIMO systems, which may hinder the training and 
fine-tuning of the network. Therefore, it is essential to develop efficient localization algorithms.

Mamba~\cite{Mamba2023}, a SSM(state-space models)-based model, provide a
promising direction for this purpose. It performs recurrent state updates with
input-dependent parameters while maintaining linear complexity with respect to
sequence length. This characteristic is attractive for indoor
localization because radio measurements can be represented as short sequences of
channel features. However, the original Mamba architecture was not designed for
Massive MIMO localization. To make it better suited for this task, in this paper, we explored 
Mamba-style architectures and simplified the original Mamba architecture to our Slim-Mamba model. 

Recent studies have further shown that Mamba-style models require dedicated
hardware-aware optimization when deployed on FPGA. Existing Mamba accelerators
have explored computation reordering, accurate quantization and pipelined execution dataflow, 
and hybrid dataflows~\cite{LightMamba2025,FastMamba2025,HCSA2025}. These works suggest
that the efficiency of Mamba acceleration depends on reusable hardware structures that maximize 
resource sharing, efficient pipelined dataflows that improve processing throughput 
and reduce inter-stage stalls, and quantization strategies that balance computational efficiency with 
inference accuracy. Motivated by this
observation, this paper explored hardware realization of Slim-Mamba for
Massive MIMO indoor localization. The proposed reconfigurable accelerator reuses computation 
resources across different projection layers, supports configurable inter-block
execution, applies a unified and configurable requantization module for different layers, and
uses fine-grained FIFO-based pipelining to reduce waiting time between projection,
nonlinear activation, state update, and output gating. 


\section{Mamba-Based Localization Models Exploration}
\label{sec:algorithm}

\section{Quantization Workflow and Hardware Implementation}
\label{sec:quant_hw}

The Slim-Mamba accelerator is implemented with fixed-point arithmetic to reduce
memory cost, DSP usage, and data-movement overhead on FPGA. The rounding, scaling, 
and fixed-point data format must also be considered to make sure the accuracy is close 
between the software model and the RTL datapath. For this reason, we use a three-level 
quantization tool, moving from a floating-point Python reference, to a bit-true C++ 
fixed-point model, and finally to RTL implementation. The Python model is used as the 
floating-point golden reference and supports fast
evaluation of different bit-widths and Q formats. The C++ fixed-point model then reproduces 
the integer behavior expected in hardware, which defines the bit-true
semantics of each operator and serves as the reference for RTL verification.

As shown in Fig.~\ref{fig:quant_tool}, by comparing the MAE between the floating-point and C++ 
models, we can get wether the 
quantization strategy is effective; by comparing the C++ and RTL outputs, we can verify whether 
the hardware mapping is consistent with the software model. This process connects the numerical 
analysis and the accelerator datapath, making the fixed-point Slim-Mamba accelerator both 
hardware-efficient and numerically consistent with the deployment-oriented software model.

\begin{figure}[!t]
  \centering
  \includesvg[width=0.85\linewidth]{TCASfig/quant_tool}
  \caption{The quantization tools: from floating-point evaluation to bit-true fixed-point modelling and RTL mapping.}
  \label{fig:quant_tool}
\end{figure}

\subsection{Quantization Workflow}

Fig.~\ref{fig:quantization_workflow} shows the proposed quantization workflow. It starts 
with a precision
sweep to determine the required precision for different computation stages. A baseline
error profiling step is then performed using a unified requantization
configuration. If the final output error is
too large, per-layer error analysis is used to locate where the mismatch grows. Based on
this analysis, the layer-wise quantization strategy is refined and then mapped to the
RTL quantization modules and memory initialization files..
A fake-quantization precision sweep is first inserted into the software model to estimate 
the accuracy impact of different quantization choices. 
In this step, quantize and
dequantize operations are inserted in the software model to emulate low-precision
effects, while the evaluation is still performed before RTL implementation. As shown in
Table~\ref{tab:precision_sweep}, the all-int16 configuration achieves a localization
error close to the fake-quantization reference. The sweep also shows that the
\texttt{dt\_proj}, \texttt{ssm\_state}, and gate paths can be reduced to int8 with
negligible loss, while \texttt{in\_proj} and \texttt{out\_proj} are more sensitive to
precision reduction. However, the accelerator uses a shared reconfigurable MAC fabric
for multiple projection stages, which will be discussed in
Section~\ref{sec:reconfigurable_hw}. To keep the array 
control, datapath width, and on-chip
memory organization simple, the final implementation adopts a unified 16-bit datapath
instead of assigning different bit-widths to different stages. This choice slightly 
sacrifices the theoretical minimum
bit-width, but it simplifies scheduling and allows the same MAC and
requantization hardware to be reused across the Slim-Mamba block.

\begin{figure}[!t]
  \centering
  \includesvg[width=0.45\linewidth]{TCASfig/quant_flow}
  \caption{Quantization workflow.}
  \label{fig:quantization_workflow}
\end{figure}

\begin{table*}[!t]
\begin{center}
\caption{Fake-Quantization Precision Sweep.}
\label{tab:precision_sweep}
\footnotesize
\setlength{\tabcolsep}{4pt}
\renewcommand{\arraystretch}{1.08}
\begin{tabular}{|p{0.17\textwidth}|p{0.50\textwidth}|>{\centering\arraybackslash}p{0.23\textwidth}|}
\hline
\textbf{Configuration} & \textbf{Precision override} & \textbf{Mean localization error (m)} \\
\hline
Fake-quant reference & Software fake-quant baseline & 0.13695 \\
\hline
All-int16 & \texttt{ssm\_state=int16}, \texttt{gate=int16} & 0.13520 \\
\hline
Mixed candidate & \texttt{dt\_proj=int8}, \texttt{ssm\_state=int8}, \texttt{gate=int8} & 0.13519 \\
\hline
\texttt{in\_proj} int8 & \texttt{in\_proj=int8}, \texttt{ssm\_state=int16}, \texttt{gate=int16} & 0.76520 \\
\hline
\texttt{out\_proj} int8 & \texttt{out\_proj=int8}, \texttt{ssm\_state=int16}, \texttt{gate=int16} & 0.24352 \\
\hline
\end{tabular}
\end{center}
\end{table*}
 



\subsection{Layer-Wise Quantization Strategy}

A unified fixed-scale baseline is first evaluated, where the main projection
outputs and selective-scan state are represented in Q8.8 and use the same scale which is 
derived from the range of int16, \(s=2^{-8}=1/256\). 
With this baseline, the
hardware fixed-point output differs substantially from the floating-point
reference, giving an output mean absolute error (MAE) of 0.17987 and a mean
localization error of 0.35543~m. To locate the dominant error source, the
Slim-Mamba block is decomposed into its main computation stages and the numerical
mismatch is measured after each stage, as presented in Table~\ref{tab:stage_quant_error}.

The stage-level analysis shows that the largest error increase occurs in the
selective-scan state update. This is expected because the state is updated
recurrently and its dynamic range varies across channels and time steps. A single
fixed Q8.8 scale is therefore insufficient to represent both small and large state
values. To reduce this error, the selective-scan state is
changed to a dynamic-scale representation. The runtime scale is derived from the
state range and used when storing and restoring the recurrent state. This reduces
the selective-scan MAE from 0.127090 to 0.085484 and also improves the following
gating and output-projection stages. The earlier stages remain unchanged because
the dynamic scaling is introduced only at the selective-scan state path.

\begin{table}[!t]
\begin{center}
\caption{Stage-Level Quantization Error With Dynamic Scaling Applied to the State Path.}
\label{tab:stage_quant_error}
\footnotesize
\setlength{\tabcolsep}{6pt}
\begin{tabular}{|l|c|c|}
\hline
\textbf{Stage} & \textbf{Fixed-scale MAE} & \textbf{Dynamic-scale MAE} \\
\hline
\texttt{norm} & 0.004498 & 0.004498 \\
\hline
\texttt{inproj\_u} & 0.010917 & 0.010917 \\
\hline
\texttt{inproj\_z} & 0.011069 & 0.011069 \\
\hline
\texttt{sigmoid\_u} & 0.006183 & 0.006183 \\
\hline
\texttt{silu\_gate} & 0.005931 & 0.005931 \\
\hline
\texttt{dtproj} & 0.010139 & 0.010139 \\
\hline
\texttt{dt\_sigmoid} & 0.002118 & 0.002118 \\
\hline
\textbf{\texttt{selective\_scan}} & \textbf{0.127090} & \textbf{0.085484} \\
\hline
\texttt{ewm\_gating} & 0.033001 & 0.029827 \\
\hline
\texttt{outproj} & 0.032225 & 0.031243 \\
\hline
\end{tabular}
\end{center}
\end{table}

The final quantization strategy is summarized in
Table~\ref{tab:quant_strategy_tcas}. Round-to-nearest-even (RNE) is used in the
projection quantization and state-conversion paths because these paths contain
repeated MAC accumulation and recurrent updates, where biased rounding errors can
accumulate across layers. The sigmoid LUT is generated offline using nearest
rounding (NR), while the runtime sigmoid path only performs clamping and table
lookup. The sigmoid output \texttt{dt\_lambda} is represented in Q0.16 and is
converted to the Q1.15 state-update coefficient $\lambda$ by a one-bit right
shift. Since $1.0$ is represented as $2^{15}$ in Q1.15, the complement
$(1-\lambda)$ can be computed directly as $2^{15}-q_{\lambda}$, avoiding an
extra quantization step.
% TODO: After Section II is written, cite the Slim-Mamba state-update equation here,
% e.g., Eq.~\eqref{eq:slim_mamba_state_update}, to justify the use of both
% $\lambda$ and $(1-\lambda)$ in the selective-scan datapath.

\begin{table}[!t]
\begin{center}
\caption{Final Fixed-Point Strategy for the Slim-Mamba Accelerator.}
\label{tab:quant_strategy_tcas}
\footnotesize
\setlength{\tabcolsep}{3pt}
\renewcommand{\arraystretch}{1.08}
\begin{tabular}{|p{0.30\linewidth}|>{\centering\arraybackslash}p{0.10\linewidth}|>{\centering\arraybackslash}p{0.20\linewidth}|>{\centering\arraybackslash}p{0.23\linewidth}|}
\hline
\textbf{Stage} & \textbf{Format} & \textbf{Scale} & \textbf{Rounding} \\
\hline
RMSNorm output       & Q8.8  & $2^{-8}$  & RNE \\
\hline
RMSNorm reciprocal   & Q30   & $2^{-30}$ & NR  \\
\hline
\texttt{in\_proj}    & Q8.8  & $2^{-8}$  & RNE \\
\hline
\texttt{dt\_proj}    & Q8.8  & $2^{-8}$  & RNE \\
\hline
\texttt{out\_proj}   & Q8.8  & $2^{-8}$  & RNE \\
\hline
Gate path            & Q8.8  & $2^{-8}$  & RNE \\
\hline
Sigmoid LUT          & Q0.16 & $2^{-16}$ & NR  \\
\hline
\texttt{dt\_lambda}  & Q0.16 & $2^{-16}$ & NR  \\
\hline
State coefficient $\lambda$ & Q1.15 & $2^{-15}$ & shift \\
\hline
Selective-scan state & int16 & dynamic & RNE \\
\hline
\end{tabular}
\end{center}
\end{table}

The improvement is also reflected in the final end-to-end evaluation in
Table~\ref{tab:quant_final_eval}. The early fixed-scale baseline has a
large output mismatch and a much higher localization error. With dynamic scaling, 
the end-to-end hardware fixed-point model achieves an
output MAE of 0.02081 with respect to the floating-point output and a mean
localization error of 0.14103~m, which is close to the
fake-quantization reference of 0.13695 m, showing that
the selected Q formats and dynamic state scaling preserve the localization
accuracy while enabling fixed-point FPGA execution.

\begin{table}[!t]
\begin{center}
\caption{End-to-End Quantization Evaluation.}
\label{tab:quant_final_eval}
\scriptsize
\setlength{\tabcolsep}{2pt}
\renewcommand{\arraystretch}{1.08}
\begin{tabular}{|p{0.18\linewidth}|>{\centering\arraybackslash}p{0.13\linewidth}|>{\centering\arraybackslash}p{0.25\linewidth}|>{\centering\arraybackslash}p{0.30\linewidth}|}
\hline
\textbf{Evaluation} & \textbf{Samples} & \textbf{HW vs. Float MAE} &
\textbf{Mean Localization Error (m)} \\
\hline
fixed scale & 32 & 0.17987 & 0.35543 \\
\hline
dynamic scale & 34,505 & 0.02081 & 0.14103 \\
\hline
Fake-quant ref. & 34,505 & -- & 0.13695 \\
\hline
\end{tabular}
\end{center}
\end{table}

\subsection{Configurable Requantization Engine}

Instead of instantiating a separate quantization unit for each layer, the
accelerator uses a shared configurable requantization engine. This module is used
after MAC accumulation, Q-format conversion, state update, scale multiplication,
and residual addition. As shown in Fig.~\ref{fig:requant_engine}, the input value
is first sign-extended or zero-extended according to the input mode. It is then
multiplied by either a fixed scale or a runtime dynamic scale. The product is
shifted to the target binary point, rounded according to the selected rounding
mode, and finally saturated or wrapped according to the output mode.

\begin{figure*}[!t]
\centering
\includesvg[width=0.78\textwidth]{TCASfig/quant_engine}
\caption{Configurable fixed-point requantization engine reused across
projection, state-update, and residual paths.}
\label{fig:requant_engine}
\end{figure*}

This shared structure reduces duplicated RTL and
keeps the rounding and allows the layer-wise quantization policy obtained from software
analysis to be adjusted easily to hardware control parameters. Different
projection stages can therefore reuse the same datapath while selecting different
scale memories, sign modes, and output formats.

The dynamic-scale state path requires additional alignment logic. Since the
selective-scan datapath is stream based, each token, state address, update
coefficient, and runtime scale must arrive at the state-update unit in the same
order. A ping-pong scale buffer is used for this purpose, see
Fig.~\ref{fig:ping_pong}.
While one bank receives
the incoming token and its generated runtime scale, the other bank provides an
already aligned token-scale pair to the downstream state-update pipeline. The
read and write banks switch between consecutive tokens, which hides the
scale-generation latency and preserves streaming order. In this way, dynamic
scaling can be inserted into the recurrent state path without breaking the
fine-grained FIFO-based pipeline.

\begin{figure}[!t]
\centering
\includesvg[width=0.95\linewidth]{TCASfig/ping-pong}
\caption{Ping-pong scale-buffer organization for aligning token, state address,
update coefficient, and runtime scale streams.}
\label{fig:ping_pong}
\end{figure}


\section{Reconfigurable Slim-Mamba Hardware Design and Pipelining}
\label{sec:reconfigurable_hw}

This section presents the reconfigurable hardware architecture of the proposed
Slim-Mamba accelerator. The term ``reconfigurable''  refers to:
\begin{itemize}
\item the same hardware chain can select different inter-block execution modes;
\item the same MAC fabric can be scheduled for different Slim-Mamba projection operators; 
\item the tile shape and reduction length of the fabric can be configured for different 
matrix-vector shape with multi-mode PEs supporting addition, multiplication, and addition after multiplication;
%\item the same multi-mode PEs inside the computing array that support addition, multiplication and addition after multiplication.
\item the same fixed-point datapath can use different scale, rounding, and saturation policies for different layers and channels, 
see Fig.~\ref{fig:requant_engine}.
\end{itemize}

Fig.~\ref{fig:hardware_architecture} shows the overall architecture. The design is
implemented on a Zynq MPSoC platform. The processing system (PS) manages
software control, DDR buffers, and AXI DMA transactions, while the programmable
logic (PL) implements the Slim-Mamba compute datapath. As shown in
Fig.~\ref{fig:hardware_architecture}(a), the contents of input buffer and output buffer are stored in
PS-side DDR. The input vector is transferred to the PL through the
memory-mapped-to-stream (MM2S) channel of the AXI DMA, and the output stream is
written back to DDR through the stream-to-memory-mapped (S2MM) channel. The PS
configures the accelerator using AXI-Lite control/status registers, including
input preload, computation start, busy/status polling, and output write-back
control.

\begin{figure*}[!t]
\centering
\includesvg[width=0.95\textwidth]{TCASfig/tcas_hardware_architecture}
\caption{Reconfigurable Slim-Mamba accelerator architecture: (a) board-level
PS-side memory and AXI interfaces, (b) PL shell and four-block Slim-Mamba
chain, (c) time-multiplexed Mamba block with shared projection fabric, and
(d) SSM block datapath.}
\label{fig:hardware_architecture}
\end{figure*}

\subsection{Board-Level Dataflow and Four-Block Execution}

The PL-side shell connects the AXI interfaces to the four-block Slim-Mamba
compute chain, as shown in Fig.~\ref{fig:hardware_architecture}(b). The AXI-Lite CSR
first starts the stream loader, which writes the incoming AXI-Stream data into
the input buffer of Block~0. After the preload stage is complete, the four-block
controller starts the block computation sequence. Each block contains a full set of operations 
in Slim-Mamba. Between adjacent blocks, the current block output is
combined with the residual path using saturated fixed-point addition and then
written into the next block input buffer. The output of the last block is
streamed to the DMA output path.

The first reconfigurable feature appears at the block boundary. The four-block
chain supports two inter-block transfer modes through the
inter block pipeline control. When
the pipeline is enabled, the accelerator uses a streaming inter-block
transfer mode. In this mode, the output of Block~(N) does not need to wait
until the full block result has been cached. Instead, a registered inter-block
write channel forwards the block output row by row or tile by tile into the
input buffer of Block~(N+1). Therefore, partial outputs from the current block
can be made available to the next block earlier, improving the opportunity for
overlap between block output generation and next-block preparation.

When the pipeline is unenabled, the accelerator uses a conservative
preload mode. The output of the current block is first accumulated or cached
completely. It is then preloaded into the input buffer of the next block, and
the next block starts only after this preload step is complete. This mode has
simpler control behavior and is useful for validation and debugging, while the
pipeline mode is more throughput-oriented. Both modes use the same four-block
hardware chain. The difference is the scheduling policy and the inter-block
write behavior, rather than the physical compute units.

\subsection{Reconfigurable Projection Fabric}

Inside each Slim-Mamba block, the main MAC-intensive operators are the input
projection, the update-gate projection in the SSM path, and the output
projection. Although these operators use different input buffers, weight banks,
and output destinations, they can all be represented as matrix-vector
multiplication:
\begin{equation}
\mathbf{y}=W\mathbf{x}+\mathbf{b}.
\label{eq:projection_mvm}
\end{equation}
The accelerator therefore avoids instantiating independent MAC arrays for
these projection layers. Instead, it uses
one shared $4\times4\times4$ MAC fabric, as illustrated in
Fig.~\ref{fig:hardware_architecture}(c).

Each projection operator has its own scheduler. The
in projection scheduler reads the block input or hidden representation and
the input-projection weights, and converts them into activation and weight tiles.
The dt projection scheduler reads the SSM input cache and the update-gate
projection weights, and generates the corresponding tile stream. The
out projection scheduler reads the gated SSM output and the output-projection
weights. Although these three schedulers access different buffers and write to
different destinations, they expose the same fabric interface. A MAC fabric
manager arbitrates among the schedulers and configures the shared fabric to
serve one projection role at a time.

This is the second reconfigurable feature of the accelerator. The physical MAC
fabric remains unchanged, but its operator role is changed by scheduler
selection, buffer selection, weight-bank selection, and output write-back
configuration. In this way, the same MAC fabric is reused across the projection
operators of the Slim-Mamba block. This reduces DSP duplication and avoids the
routing and memory pressure that would arise from placing separate projection
arrays in the same block.

The third reconfigurable feature is the configurable tile shape and reduction
behavior of the same fabric. The fabric interface includes control parameters
such as \texttt{mode}, \texttt{col\_blocks\_cfg}, and \texttt{reduce\_rows}.
The \texttt{mode} signal selects the current PE operation, such as MAC,
EWM(element-wise multiplication), or EWA(element-wise addition). In the current
accelerator, the shared fabric is mainly used in MAC mode for projection
operators, while lightweight vector modules handle the fine-grained SSM
element-wise path. However, the mode interface keeps the PE structure
configurable and allows the same fabric abstraction to support different operator
types.

The \texttt{col\_blocks\_cfg} parameter controls how many column blocks must be
accumulated for one output tile. This is necessary because different projection
layers have different input dimensions and therefore different reduction lengths
along the $K$ dimension. Rather than fixing the reduction length inside the PE
array, each scheduler provides the required column-block count for its own
projection workload. The \texttt{reduce\_rows} signal controls whether the array
outputs are reduced across rows to generate the final vector output. Together,
these parameters allow the fixed $4\times4\times4$ physical array to support
different logical matrix shapes, reduction lengths, and output forms.

The $4\times4\times4$ fabric contains four $4\times4$ sub-arrays. Each
sub-array receives a four-element activation slice and a $4\times4$ weight
block. The four sub-arrays operate with a staggered schedule and their partial
results are combined through a reduction tree. Figure \ref{fig:array_organization} 
illustrates the array organization for dt projection in SSM path. The B in the table
stands for the current input, which is a $4\times4$ block. The weight matrix below is
part of the weight mastrix for the dt projection, and each square means a $4\times4$ block. 
The four sub-arrays work in pipeline. At cycle 19, the first 4 rows of the weight matrix is
completely consumed, and at cycle 22, the first output tile is produced after the reduction tree.

Because the main calculation in the 
projection operators is matrix-vector multiplication, compared with systolic 
array, our $4\times4\times4$ arrays organization reduces global routing and
data-supply pressure for the MVM-dominated Slim-Mamba workload. Compared with a
$64\times1$ GEMV engine, it consumes the same amount of hardware resources but provides
earlier tile availability: still take SSM as example, which
performs MVM with the matrix shape of (256,256) and vector shape of (256,1). With the 
four $4\times4$ sub-arrays working in pipeline, the first output vector can be produced 
after 22 cycles rather than waiting for a full
vector-reduction interval in a $64\times1$ GEMV engine, which is 256 cycles. 
This earlier tile output is important because the
following SSM nonlinear and state-update stages can start before the whole
projection vector is completed.

\begin{figure*}[!t]
\centering
\includesvg[width=0.75\textwidth]{TCASfig/small_parallel_pipeline}
\caption{$4\times4\times4$ arrays organization.}
\label{fig:array_organization}
\end{figure*}


\subsection{Fine-Grained FIFO-Based Pipeline}

As we memtioned before, the PEs inside the $4\times4\times4$ MAC arrays can be configured 
for MVM, EWM, EWA, so the trade-off between hardware reuse and latency needs to be considered. 
A coarse implementation of the SSM path would execute update-gate projection,
sigmoid, state update, SiLU, and output gating sequentially. This maximizes
resource reuse, but it forces each stage to wait until the previous stage has
produced a full vector. The proposed accelerator instead uses a fine-grained
FIFO-based pipeline. Once the MAC fabric produces and reduces an output tile,
the tile is forwarded through FIFO interfaces to the bias-add, sigmoid, and
state-update stages. The state-update unit starts as soon as the corresponding
tile, update coefficient, state value, and scale are available. It does not wait
for the full projection vector.

The state update follows the Slim-Mamba recurrence
\begin{equation}
h_t = \lambda_t \odot h_{t-1}
+(1-\lambda_t)\odot x_t^{(s)} ,
\label{eq:slim_mamba_state_update}
\end{equation}
where $\lambda_t$ is generated by the update-gate projection and sigmoid
unit. The updated state is then multiplied with the SiLU gate output. FIFOs are
used to align the projection output, sigmoid value, gate path, state read result,
and dynamic scale. This alignment is necessary because different sub-stages have
different local latencies.

For an $N$-tile SSM workload, the coarse-grained schedule can be abstracted as
\begin{equation}
T_{\mathrm{coarse}} = 29N ,
\label{eq:coarse_latency}
\end{equation}
where projection, nonlinear activation, state update, and output gating are
mostly serialized. With fine-grained FIFO streaming, the schedule becomes
\begin{equation}
T_{\mathrm{fine}} = 22N + 7 .
\label{eq:fine_latency}
\end{equation}
For $N=64$, this reduces the SSM-path latency from 1856 cycles to 1415 cycles,
corresponding to a 23.8\% reduction. The additional cost is about 12 DSPs per
SSM compute path, which is acceptable compared with the latency reduction. 
~Figure~\ref{fig:fine_grained_pipeline} shows the comparison between 
coarse-grained and fine-grained scheduling. Fig.~\ref{fig:hardware_architecture}(d) 
shows the final fine-grained SSM block datapath, where projection,
bias addition, sigmoid, state update, SiLU gating, and output generation are
connected through FIFO-based interfaces.

\begin{figure}[!t]
\centering
\includesvg[width=0.95\linewidth]{TCASfig/pipeline_scheme}
\caption{Comparison between coarse-grained and fine-grained scheduling.}
\label{fig:fine_grained_pipeline}
\end{figure}

Overall, the accelerator combines four forms of architecture-level
reconfigurability. The block chain can switch between streaming and preload
execution. The shared MAC fabric can be configured by schedulers to execute
different projection operators. The tile and reduction behavior of the fabric
can adapt to different projection shapes. The quantization path can select
different scale, rounding, and saturation policies for different layers and
channels. These mechanisms allow the Slim-Mamba accelerator to reuse most of its
hardware resources while preserving a fine-grained streaming path through the
latency-sensitive SSM computation.

\subsection{Nonlinear Layer Implementation}

From the algorithm analyse in section \ref{sec:algorithm}, the main nonlinear operators insede 
our Slim-Mamba model are RMSNorm, sigmoid, and SiLU.

RMSNorm is placed before the projection path to stabilize the dynamic range of
the input features. For an input vector
$\mathbf{x}=[x_0,x_1,\ldots,x_{D-1}]$, the RMS value is
\begin{equation}
\mathrm{RMS}(\mathbf{x})=
\sqrt{\frac{1}{D}\sum_{i=0}^{D-1}x_i^2+\epsilon},
\label{eq:rms_value}
\end{equation}
and the normalized output is computed as
\begin{equation}
\hat{x}_i =
x_i \cdot
\widehat{\frac{1}{\mathrm{RMS}(\mathbf{x})}} .
\label{eq:rmsnorm_output}
\end{equation}
The square-root operation is implemented with a trial-quotient digit-recurrence
procedure, which uses shift, comparison, and subtraction operations instead of a
general floating-point square-root unit, as shown in Figure \ref{fig:sqrt}. The final 
division is replaced by reciprocal generation followed by multiplication. A Q30 reciprocal 
is used so that the shared normalization scale remains accurate for all channels of one
token.

\begin{figure}[!t]
\centering
\includesvg[width=0.75\linewidth]{TCASfig/rms_approximation_flow}
\caption{RMSNorm approximation flow: a trial-quotient digit-recurrence
procedure.}
\label{fig:sqrt}
\end{figure}

The sigmoid function is used to generate the state-update coefficient. It is
implemented with a lookup table:
\begin{equation}
\sigma(x)=\frac{1}{1+e^{-x}}.
\label{eq:sigmoid}
\end{equation}
The input is clamped to $[-4,+4)$, and a 2048-entry LUT stores the
corresponding Q0.16 sigmoid values. This avoids evaluating the exponential
function in hardware. SiLU reuses the same sigmoid path:
\begin{equation}
\mathrm{SiLU}(x)=x\cdot\sigma(x).
\label{eq:silu}
\end{equation}
The original input is delayed through a FIFO while the sigmoid value is produced.
After alignment, an element-wise multiplier generates the SiLU output. Therefore,
sigmoid and SiLU share the clamp and LUT logic, while SiLU only adds FIFO
alignment and vector multiplication.



\section{Results}

\section{Conlusions}

\section*{Acknowledgments}
This should be a simple paragraph before the References to thank those individuals and institutions who have supported your work on this article.



{\appendix[Proof of the Zonklar Equations]
Use $\backslash${\tt{appendix}} if you have a single appendix:
Do not use $\backslash${\tt{section}} anymore after $\backslash${\tt{appendix}}, only $\backslash${\tt{section*}}.
If you have multiple appendixes use $\backslash${\tt{appendices}} then use $\backslash${\tt{section}} to start each appendix.
You must declare a $\backslash${\tt{section}} before using any $\backslash${\tt{subsection}} or using $\backslash${\tt{label}} ($\backslash${\tt{appendices}} by itself
 starts a section numbered zero.)}



%{\appendices
%\section*{Proof of the First Zonklar Equation}
%Appendix one text goes here.
% You can choose not to have a title for an appendix if you want by leaving the argument blank
%\section*{Proof of the Second Zonklar Equation}
%Appendix two text goes here.}



\bibliographystyle{IEEEtran}
\bibliography{thesisbib}


\newpage

\section{Biography Section}
If you have an EPS/PDF photo (graphicx package needed), extra braces are
 needed around the contents of the optional argument to biography to prevent
 the LaTeX parser from getting confused when it sees the complicated
 $\backslash${\tt{includegraphics}} command within an optional argument. (You can create
 your own custom macro containing the $\backslash${\tt{includegraphics}} command to make things
 simpler here.)
 
\vspace{11pt}

\bf{If you include a photo:}\vspace{-33pt}
\begin{IEEEbiography}[{\includegraphics[width=1in,height=1.25in,clip,keepaspectratio]{fig1}}]{Michael Shell}
Use $\backslash${\tt{begin\{IEEEbiography\}}} and then for the 1st argument use $\backslash${\tt{includegraphics}} to declare and link the author photo.
Use the author name as the 3rd argument followed by the biography text.
\end{IEEEbiography}

\vspace{11pt}

\bf{If you will not include a photo:}\vspace{-33pt}
\begin{IEEEbiographynophoto}{John Doe}
Use $\backslash${\tt{begin\{IEEEbiographynophoto\}}} and the author name as the argument followed by the biography text.
\end{IEEEbiographynophoto}




\vfill

\end{document}
